% Imporditud: tools/import_eksperimendid.py
% Allikas: Eesti füüsikaolümpiaadi eksperimendi ülesannete
%   kogu (Rael Kalda, 2023), ülesanne Ü135, lahendus L135.
\setAuthor{EFO žürii}
\setRound{lõppvoor}
\setYear{2008}
\setNumber{G E1}
\setDifficulty{4}
\setTopic{elektriahelad}
\setVahendid{patarei (sisetakistus on tühine), tuntud takisti (R = 4,7 k$\Omega \pm$ 5\%), voltmeeter, ühendusjuhtmed}
\setPunktid{10}
\setAllikas{Eksperimendi kogu Ü135, lahendus lk 103}
\setLahendusseis{toores}

\prob{Ampermeeter}
Voltmeetrit tahetakse kasutada milliampermeetrina. Joonistage selleks sobiv skeem ning määrake, mitmele milliamprile vastab sel juhul voltmeetri näidu üks volt. Hinnake mõõteviga.

\hint

\solu
Voltmeetri milliampermeetrina kasutamiseks tuleb ta lEütliIta=daUp(arR1al+lee1rl)s.eKltutiakmisetiiega R. Mõõdame voltmeetri sisetakistuse r (vt allpool).

tehtud ampermeetri klemmidele tuleb vool I ning voltmeetri lugem on U siis keh-

UU tib seos I = + . Seega on 1 V puhul voolutugevus läbi meie ampermeetri R( r ) klemmide 103 $\cdot$ + [mA/V]. Rr

Niisiis peame leidma voltmeetri sisetakistuse. Selleks rakendame patarei pinge es-

malt otse voltmeetrile, lugem $U_1$ = E annab patarei elektromotoorjõu; seejärel

rakendame jadamisi takistiga R, lugem $U_2$ = rE/(r + R) = rU1/(r + R). Seega

1 + R/r = $U_1$/$U_2$, millest r = R $U_2$ . $U_1$ $-$ $U_2$

Asendades see tulemus eelpool saadud valemisse näeme, et 1 V-se näidu puhul on

voolutugevus

103 $\cdot$ $U_1$ [mA/V]. U2R
\probend
